\chapter{Unix Shell}
\begin{center}
    \includegraphics[width=120px]{img/term.pdf} \\
    \textbf{\href{http://en.wikipedia.org/wiki/Unix_shell}{Unix Shell}}
\end{center}

\section{Dateisystem}
Das Dateisystem in Linux besteht aus \emph{einem} Baum.
Dies äußert sich im so genannten Dateipfad, der den Weg von der Wurzel \texttt{/} zur eigentlichen Datei oder dem Ordner beschreibt.
Nach jedem Knoten, das heißt nach jedem Überordner setzt man wieder ein Slash bis man beim Dateinamen angekommen ist.

Es gibt immer ein aktuelles Verzeichnis, in dem man sich befindet; dieses wird 'working directory' genannt.
Will man aus dem working directory in ein anderes Verzeichnis so kann man dies durch eine relative oder absolute Pfadangabe tuen.
Bei einem absoluten Pfad geht man wieder von der Wurzel aus und beginnt ihn mit \texttt{/}.
Wählt man den relativen Pfad, der oft kürzer ist, insbesondere wenn man einfach nur in den Über- oder Unterordner möchte, so baut man die Pfadangabe ausgehend vom aktuellen Verzeichnis auf.
Bei der Pfadangabe sollte man im Übrigen folgende besondere Verzeichnisse kennen:
    \begin{itemize}
      \item \texttt{.} das aktuelle Verzeichnis (oder der aktuelle im bisherigen Pfad, \texttt{a/./}~=~\texttt{a/})
      \item \texttt{..} das Oberverzeichnis des aktuellen Verzeichnisses (\texttt{a/b/../}~=~\texttt{a/})
      \item \texttt{\textasciitilde} das Heimverzeichnis (nur am Anfang eines Pfads)
    \end{itemize}

\section{Aufbau einer Eingabe}
\texttt{\$ ls -l --all \textit{directory}\\
        \textit{output}\\
        \$}
\begin{center}
  \begin{tabular}{>{\tt}l l}
    \toprule
    \$                 & Prompt       \\
    ls                 & Befehl       \\
    -l                 & kurze Option \\
    --all              & lange Option \\
    \textit{directory} & Argument     \\
    \textit{output}    & Ausgabe      \\
    \bottomrule
  \end{tabular}
\end{center}
\begin{itemize}
  \item \texttt{\$} ist nur ein Beispiel für einen Prompt, häufig wird das aktuelle Verzeichnis und/oder andere Informationen angezeigt
  \item kurze Optionen können zusammengefasst werden (\texttt{ls~-la} = \texttt{ls -l -a} = \texttt{ls -l --all})
  \item die Reihenfolge der Optionen ist egal
  \item meistens werden mehrere Argumente (z.B. Dateien) akzeptiert
\end{itemize}

\section{Befehle}
\subsection{\texttt{man}}
\begin{itemize}
  \item \texttt{man \textit{topic}} für manual: zeigt die Hilfe für ein Programm
  \item Beispiel: \texttt{man man}
\end{itemize}

\subsection{\texttt{pwd}, \texttt{cd}}
\begin{itemize}
  \item \texttt{pwd} für print working directory: zeigt das aktuelle Verzeichnis
  \item \texttt{cd \textit{directory}} für change directory: wechselt in das angegebene Verzeichnis
  \item Beispiel:
\begin{verbatim}
$ pwd
/home/ismo
$ cd ../../etc
$ pwd
/etc
$ cd ~
$ pwd
/home/ismo
\end{verbatim}
\end{itemize}

\subsection{ls}
\begin{itemize}
  \item \texttt{ls [\textit{directory}]} für list: zeigt den Inhalt eines Verzeichnisses an
  \item \texttt{ls -l}: zeigt mehr Informationen über Dateien und Verzeichnisse
  \item \texttt{ls -a}: zeigt auch versteckte Dateien
  \item \texttt{ls -R}: zeigt auch den Inhalt von Unterverzeichnissen
  \item alle Optionen können kombiniert werden
\end{itemize}
Beispiel:
\begin{verbatim}
$ ls
a/  b
$ ls -l
total 4.0K
drwxr-xr-x 2 ismo users 4.0K Sep 15 19:52 a/
-rw-r--r-- 1 ismo users    0 Sep 15 19:52 b
$ ls -a
./  ../  a/  b
$ ls -R
.:
a/  b

./a:
c
\end{verbatim}

\subsection{\texttt{mkdir}, \texttt{touch}}
\begin{itemize}
  \item \texttt{mkdir \textit{directory}} für make directory: erstellt ein neues Verzeichnis
  \item \texttt{mkdir -p \textit{directory}}: erstellt ein neues Verzeichnis und alle notwendigen Oberverzeichnisse
  \item \texttt{touch \textit{file}}: erstellt eine neue, leere Datei
\end{itemize}
Beispiel:
\begin{verbatim}
$ ls
$ mkdir a
$ mkdir b/c
mkdir: cannot create directory ‘b/c’: No such file or directory
$ mkdir -p b/c
$ touch b/file
$ ls -R
.:
a/  b/
~
./a:
~\\                  
./b:
c/ file

./b/c:
\end{verbatim}

\subsection{\texttt{cp}, \texttt{mv}, \texttt{rm}, \texttt{rmdir}}
\begin{itemize}
  \item \texttt{cp \textit{source} \textit{destination}} für copy: kopiert eine Datei
  \item \texttt{cp -r \textit{source} \textit{destination}}: kopiert ein Verzeichnis rekursiv
  \item das Ziel kann ein Verzeichnis oder der exakte Pfad sein
  \item \texttt{\textit{source}} können mehrere Dateien sein, der letzte Pfad zählt als \texttt{\textit{destination}}
  \item \texttt{mv \textit{source} \textit{destination}} für move: verschiebt order benennt eine Datei um
  \item Ziel kann wir bei \texttt{cp} sein
  \item \texttt{rm \textit{file}} für remove: löscht eine Datei
  \item \texttt{rm -r \textit{file}}: löscht ein Verzeichnis rekursiv
  \item \texttt{rmdir \textit{directory}} für remove directory: löscht ein leeres Verzeichnis
  \item \texttt{rm -r} kann statt \texttt{rmdir} verwendet werden
\end{itemize}
Beispiel:
\begin{verbatim}
$ ls
a
$ cp a b
$ ls
a  b
$ mv b c
$ ls
a  c
$ rm a
$ ls
c
\end{verbatim}

\subsection{\texttt{cat}, \texttt{less}, \texttt{grep}, \texttt{echo}}
\begin{itemize}
  \item \texttt{cat [\textit{file}]} für concatenate: gibt den Inhalt einer (oder mehrerer) Dateien aus
  \item \texttt{less [\textit{file}]} (besser als \texttt{more}): zeigt eine Datei in einer navigablen Form an
  \item \texttt{grep \textit{pattern} [\textit{file}]} für ???: sucht nach einem Muster
  \item \texttt{grep -i \textit{pattern} [\textit{file}]}: ignoriert Groß- und Kleinschreibung
  \item \texttt{grep -r \textit{pattern} \textit{directory}}: sucht rekursiv in allen Dateien
  \item \texttt{echo \textit{text}}: gibt einen Text aus
\end{itemize}

\section{Nützliche Shell Features}
\subsection{Tastaturkürzel}
\begin{itemize}
  \item \texttt{Ctrl-C}: beendet das laufende Programm
  \item \texttt{Ctrl-D}: EOF (end of file) eingeben, kann Programme beenden
  \item \texttt{Ctrl-L}: leert den Bildschirm
\end{itemize}

\subsection{Ein- und Ausgabe}
\begin{itemize}
  \item Programme haben einen Eingabe- und einen Ausgabestream
  \item diese können verschieden verwendet werden
  \item \texttt{\textit{command} > \textit{file}}: überschreibt eine Datei mit der Ausgabe
  \item \texttt{\textit{command} >> \textit{file}}: wie \texttt{>}, aber schreibt ans Ende der Datei, statt zu überschreiben
  \item \texttt{\textit{command} < \textit{file}}: verwendet eine Datei als Eingabe
  \item \texttt{\textit{command1} | \textit{command2}}: verwendet die Ausgabe eines Programms als Eingabe eines anderen
\end{itemize}
Beispiel:
\begin{verbatim}
$ echo "bla bla bla" > a
$ cat a
bla bla bla
$ echo "foo" >> a
$ cat a
bla bla bla
foo
$ grep foo < a
foo
$ touch abc
$ touch bcd
$ touch cde
$ ls | grep b
abc
bcd
\end{verbatim}

\subsection{Globbing (\texttt{*})}
\begin{itemize}
  \item \texttt{*} wird ersetzt durch alle passenden Dateien
  \item Beispiel:
\begin{verbatim}
$ ls .vim*
.viminfo  .vimrc

.vim:
backup/  bundle/  swap/  syntax/
\end{verbatim}
\end{itemize}

\section{Weitere nützliche Befehle}
\begin{itemize}
  \item \texttt{vim}: Ein sehr mächtiger Texteditor in der Kommandozeile
  \item \texttt{vimtutor}: Einführung in \texttt{vim}
  \item \texttt{emacs}: Noch ein Editor
  \item \texttt{nano}: Einfacher Editor in der Kommandozeile
  \item \texttt{find}, \texttt{locate}: Finden Dateien und Verzeichnisse
  \item \texttt{curl}, \texttt{wget}: Dateien herunterladen
  \item \texttt{rsync}: Synchronisierung von Verzeichnisstrukturen
  \item \texttt{awk}, \texttt{sed}: Bearbeitung von Text
  \item \texttt{ln}: Erstellt Links zwischen Dateien und Verzeichnissen
  \item \texttt{diff}: Zeigt  unterschiede zwischen Dateien
  \item \texttt{patch}: Wendet einen von \texttt{diff} erstellten Diff an
\end{itemize}

\section{Weiterführende Links}
\begin{itemize}
  \item Beginning with the Shell: \url{http://youtu.be/Sye3mu-EoTI}
  \item Learn CLI The Hard Way: \url{http://cli.learncodethehardway.org/book/}
  \item Bash Guide for Beginners: \url{http://tldp.org/LDP/Bash-Beginners-Guide/html/index.html}
  \item Advanced Bash-Scripting Guide: \url{http://tldp.org/LDP/abs/html/index.html}
  \item Zsh Workshop: \url{http://www.acm.uiuc.edu/workshops/zsh/toc.html}
  \item Zsh Manual: \url{http://zsh.sourceforge.net/Doc/Release/zsh\_toc.html}
  \item To understand the command line…: \url{http://geekblog.oneandoneis2.org/index.php/2012/09/30/to-understand-the-command-line}
  \item The Unix-Haters Handbook: \url{http://m.simson.net/ugh.pdf}
\end{itemize}
